\appendix
\chapter{Appendix}

\section{Numerical Computations: Python Code}

The following Python code reproduces the numerical values of this work
(arc lengths, torsion, curvature). It can be executed directly and
requires only the standard libraries \texttt{numpy} and \texttt{scipy}.

\begin{verbatim}
import numpy as np
from scipy import integrate

# -------------------------------------------------
# Basic functions
# -------------------------------------------------
def k(n):
    """Quadratic mean-value sequence."""
    return (2*n**2 + 3*n + 1) / 6

def theta_n(n):
    """Angular parameter for integer points on the helix."""
    return (np.pi / 3) * n + (np.pi / 4)

# Helix parameters
a = 6 / np.pi**2   # Spiral slope
b = 3 / np.pi       # Height slope

# -------------------------------------------------
# Arc length of the spiral (p-parametrisation)
# -------------------------------------------------
def arc_length_spiral(p1, p2):
    """Exact arc-length integral of the spiral."""
    integrand = lambda p: (1/np.pi) * np.sqrt(3/p + np.pi**2)
    L, _ = integrate.quad(integrand, p1, p2)
    return L

def arc_length_spiral_taylor(p1, p2):
    """Taylor approximation of the arc length."""
    return ((p2 - p1)
            + (3 / (2*np.pi**2)) * np.log(p2/p1)
            + (9 / (8*np.pi**4)) * (1/p1 - 1/p2))

# -------------------------------------------------
# Arc length of the helix (theta-parametrisation)
# -------------------------------------------------
def arc_length_helix(t1, t2):
    """Arc-length integral of the conical helix."""
    integrand = lambda t: (3/np.pi**2) * np.sqrt(4*(1+t**2) + np.pi**2)
    L, _ = integrate.quad(integrand, t1, t2)
    return L

def arc_length_parabola(n1, n2):
    """Arc length of the quadratic parabola k(n)."""
    integrand = lambda n: (1/6) * np.sqrt(36 + (4*n+3)**2)
    L, _ = integrate.quad(integrand, n1, n2)
    return L

# -------------------------------------------------
# Curvature and torsion
# -------------------------------------------------
def curvature(t):
    """Planar spiral curvature (2D approximation, cf. remark in Ch. 13)."""
    return (np.pi**2 * (t**2 + 2)) / (6 * (t**2 + 1)**1.5)

def curvature_3d(t):
    """Exact 3D curvature of the conical helix (Theorem 13.9)."""
    c1 = 3 / np.pi
    num = a * np.sqrt(c1**2 * (t**2 + 4) + a**2 * (t**2 + 2)**2)
    den = (a**2 * (1 + t**2) + c1**2)**1.5
    return num / den

def torsion(t):
    """Torsion of the conical helix."""
    num = np.pi**3 * (6 + t**2)
    den = 3 * (np.pi**2 * (4 + t**2) + 4*(2 + t**2)**2)
    return num / den

# -------------------------------------------------
# Example outputs
# -------------------------------------------------
if __name__ == "__main__":
    print("=== Difference sequence (major/minor steps) ===")
    for m in range(1, 6):
        maj = 16*m - 6
        mi  = 8*m + 1
        Um  = maj + mi
        print(f"m={m}: Major={maj}, Minor={mi}, U_m={Um}")

    print("\n=== Arc-length comparison: parabola vs helix ===")
    pairs = [(1,5),(5,7),(11,13),(17,19),(23,25),(35,37)]
    for n1, n2 in pairs:
        t1 = theta_n(n1)
        t2 = theta_n(n2)
        Lp = arc_length_parabola(n1, n2)
        Lh = arc_length_helix(t1, t2)
        dev = (Lh - Lp) / Lp * 100
        print(f"n={n1}->{n2}: L_Parabola={Lp:.6f}, L_Helix={Lh:.6f}, "
              f"Dev={dev:+.2f}%")

    print("\n=== Curvature and torsion at prime-number angles ===")
    primes = [5, 7, 11, 13, 17, 19, 23]
    for p in primes:
        t = theta_n(p)
        print(f"p={p:3d}: theta={t:.4f}, "
              f"kappa_2D={curvature(t):.6f}, "
              f"kappa_3D={curvature_3d(t):.6f}, "
              f"tau={torsion(t):.6f}")
\end{verbatim}

\section{Numerical Verifications}

The following tables supplement the above Python code with independent
comparisons that confirm the analytical formulae from the respective
main chapters.

\subsection{Torsion Formula (Theorem~\ref{thm:3d-curvature-torsion})} 

The torsion formula
\[
\tau(t) = \frac{\pi^3(6+t^2)}{3[\pi^2(4+t^2)+4(2+t^2)^2]}
\]
was verified by direct comparison with the numerical computation via the
cross product $(\gamma' \times \gamma'') \cdot \gamma''' / |\gamma' \times \gamma''|^2$.
The following values agree to at least $10$ decimal places:

\begin{table}[h]
\centering
\caption{Verification of the torsion at selected points}
\begin{tabular}{ccc}
\toprule
$t$ & $\tau_{\text{formula}}$ & $\tau_{\text{numerical}}$ \\
\midrule
$5$  & $0.100055070\ldots$ & $0.100055070\ldots$ \\
$10$ & $0.025691658\ldots$ & $0.025691658\ldots$ \\
$20$ & $0.006451662\ldots$ & $0.006451662\ldots$ \\
$47$ & $0.001169444\ldots$ & $0.001169444\ldots$ \\
\bottomrule
\end{tabular}
\end{table}

\subsection{Difference Formulae (Theorem~\ref{satz:minor_major_formeln})}

The formulae $\Delta_{\mathrm{major},m} = 16m-6$ and $\Delta_{\mathrm{minor},m} = 8m+1$
were verified by direct computation from $k(6m-1)-k(6m-5)$ and $k(6m+1)-k(6m-1)$
for $m = 1, \ldots, 20$ (complete agreement).
